\section{Discussion}

This follow-up is meant to sharpen, not broaden indiscriminately, the contribution of the existing benchmark. The earlier two-study paper already shows that governance can be evaluated under repeated adversarial strategy injection in sequential commons, that central signals can be ranked in a minimal setting, and that a narrower hybrid-versus-top-down architecture comparison remains informative under a stronger attacker. The next paper should therefore avoid retelling the full project chronology. Its purpose is to isolate the more ambitious question that the current results already point toward: which institutional arrangements remain robust as strategic pressure rises in a plural, interacting system.

That emphasis changes both the intellectual positioning and the standards of evidence. It shifts the paper toward institutional alignment and scalable oversight in multi-agent settings, while still keeping the claims empirical and benchmark-level. It also raises the bar for the Harvest evidence. A serious architecture paper should not stop at one pairwise comparison if the environment already supports a richer spectrum. It should not report a welfare cost without showing who bears it. And it should not discuss stronger attackers only implicitly when the codebase already contains a usable capability ladder.

At the same time, the follow-up should retain strict scope discipline. It should not claim to study all forms of governance, all possible social environments, or all relevant notions of alignment. It should not treat hybrid governance as universally dominant, and it should not promote Harvest LLM evidence into the main result chain before the reliability gates are cleared. The benchmark remains intentionally stylized, and its value lies in identification and controlled comparison rather than in direct realism.

If the next experiment cycle succeeds, the resulting paper will make a cleaner claim than the current two-study draft. Fishery Commons will establish that central interventions can be ranked under adversarial turnover. Harvest Commons will then show, more directly, how governance architecture responds to rising capability pressure and how the benefits and costs of robust governance are distributed. That is a stronger institutional-alignment story than the current paper can responsibly tell, and it is achievable without discarding the evidence already in the repository.
